\chapter{Backend abstraction}
\label{ch:backend}

\begin{aside}
\maya{} mostly talks to a real terminal. Sometimes it talks to a
file, a string, or a fake terminal in a unit test. The backend
abstraction makes the renderer not care.
\end{aside}

\section{The backend interface}

\begin{lstlisting}
class Backend {
public:
    virtual void write(std::string_view bytes) = 0;
    virtual Size size() const = 0;
    virtual Capabilities caps() const = 0;
    virtual std::optional<Event> poll(Duration) = 0;
    virtual ~Backend() = default;
};
\end{lstlisting}

Four methods. The renderer emits via \code{write}; layout queries
\code{size}; the colour and protocol selection consults
\code{caps}; the event loop blocks on \code{poll}.

\section{Real terminal backend}

The default backend wraps fd 0 (stdin) and fd 1 (stdout):

\begin{itemize}[leftmargin=*]
  \item \code{write} is a thin wrapper over \code{::write(1, ...)}.
  \item \code{size} calls \code{ioctl(TIOCGWINSZ)}.
  \item \code{caps} comes from environment detection and
    runtime queries.
  \item \code{poll} uses \code{poll(2)} on stdin.
\end{itemize}

\section{String backend}

For testing, a backend that writes to an in-memory string and
fakes input:

\begin{lstlisting}
class StringBackend : public Backend {
    std::string output;
    std::queue<Event> events;
public:
    void write(std::string_view b) override {
        output.append(b);
    }
    void inject(Event e) { events.push(e); }
    // ...
};
\end{lstlisting}

Tests construct a \code{StringBackend}, inject simulated key
events, drive the renderer, and assert on the output string.

\section{File backend}

A backend that writes ANSI to a file. Useful for recording
sessions for later replay:

\begin{lstlisting}
auto rec = std::make_unique<FileBackend>("session.ansi");
\end{lstlisting}

Playback is \code{cat session.ansi}; the terminal interprets
the recorded ANSI. Tools like asciinema do something similar
with extra timing data.

\section{TTY detection}

If stdout is not a terminal (\code{isatty} returns false), the
default backend degrades:

\begin{itemize}[leftmargin=*]
  \item \code{caps} returns minimal (no colour, no mouse).
  \item \code{size} returns \code{80x24} (a guess).
  \item Mode switches (alt-screen, mouse) are no-ops.
\end{itemize}

This lets a TUI tool work in pipes:
\code{myapp | grep something} produces plain text output.

\section{Multiple backends in one process}

A complex app might want to render to multiple terminals at
once (a remote-debug feature: ``mirror my screen to a
support engineer''). Each terminal gets its own renderer with
its own backend, sharing the element tree.

\maya{} supports this via composition: construct two
\code{Renderer} objects with different backends. The element
tree is owned by the application; both renderers walk it.

\section{Networked backends}

A backend that writes to a socket can drive a remote
terminal. This is what SSH does already (ANSI bytes over the
wire), but a custom protocol could compress better or carry
metadata.

\maya{} does not ship a network backend. The interface
supports it; implementations are application-specific.

\section{Headless rendering}

For documentation generation, render frames to an in-memory
canvas and dump the canvas as text or HTML. The string
backend collects the ANSI; a separate post-processor parses
it back into a grid.

\maya{}'s test suite uses this to verify output without a
real terminal.

\section{Backend capabilities cache}

\code{caps()} is called many times per frame (every SGR
emit checks colour mode, for instance). The default backend
caches the result; the cache is invalidated only on
explicit refresh.

\section{Recap}

\begin{recap}
\begin{itemize}[leftmargin=*]
  \item Backend = write + size + caps + poll.
  \item Real terminal, string, file, network: same interface.
  \item Pipe detection produces a degraded backend.
  \item Multiple backends per process for mirroring.
  \item Headless backends enable snapshot testing.
\end{itemize}
\end{recap}

\section*{Workshop}

\begin{enumerate}[leftmargin=*]
  \item Wire a \code{StringBackend} into a unit test.
    Inject keys and assert on the output.
  \item Build a file-backend variant that timestamps each
    write; replay with delays to simulate the original
    timing.
  \item Detect pipe vs terminal and run your app in both;
    confirm the degraded path produces useful output.
\end{enumerate}
